\section*{अध्याय \ref{ch:b1:fractions} --- परिमेय भिन्नें}

\begin{solution}{exo:b1:fractions:1}
$\dfrac{1}{X^2 - 1}$: सरल \omterm{def:b1:fractions:field}{ध्रुव} $\pm 1$; आवरण विधि: $\dfrac{1/2}{X-1} -
\dfrac{1/2}{X+1}$।

$\dfrac{X}{X^2 - 3X + 2} = \dfrac{X}{(X-1)(X-2)}$: आवरण विधि $1$ पर
$\frac{1}{1-2} = -1$ और $2$ पर $\frac{2}{2-1} = 2$ देती है: $\dfrac{-1}{X-1}
+ \dfrac{2}{X-2}$।

$\dfrac{X^2+1}{X(X-1)}$: घात $0$ है, अतः पूर्णांक भाग है: भाग देने पर $X^2 +
1 = (X^2 - X) + (X + 1)$, अतः $F = 1 + \frac{X+1}{X(X-1)}$। शेष पर आवरण
विधि: $0$ पर $\frac{1}{-1} = -1$, $1$ पर $\frac{2}{1} = 2$:
\[
F = 1 - \frac{1}{X} + \frac{2}{X-1} .
\]
\end{solution}

\begin{solution}{exo:b1:fractions:2}
$\dfrac{1}{X(X^2+1)}$: आकार $\frac aX + \frac{bX + c}{X^2 + 1}$। $0$ पर आवरण
विधि: $a = 1$। $xF$ की सीमा: $0 = a + b$, अतः $b = -1$। $X = 1$ पर मान:
$\frac12 = 1 + \frac{c - 1}{2}$, अतः $c = 0$:
\[
\frac{1}{X(X^2+1)} = \frac 1X - \frac{X}{X^2+1} .
\]

$\dfrac{X^3}{X^2+X+1}$: भाग: $X^3 = (X^2+X+1)(X - 1) + 1$, अतः
\[
\frac{X^3}{X^2+X+1} = X - 1 + \frac{1}{X^2 + X + 1} ,
\]
जो पहले से वास्तविक वियोजित रूप में है (द्विघात का विविक्तकर ऋणात्मक है)।
\end{solution}

\begin{solution}{exo:b1:fractions:3}
$\dfrac{1}{X^2(X-1)}$: आकार $\frac{a}{X^2} + \frac bX + \frac{c}{X-1}$।
द्विक \omterm{def:b1:fractions:field}{ध्रुव} $0$ पर आवरण विधि: $a = \bigl[\frac{1}{X-1}\bigr]_{0} = -1$। $1$
पर आवरण विधि: $c = 1$। $xF$ की सीमा: $0 = b + c$, अतः $b = -1$:
\[
\frac{1}{X^2(X-1)} = -\frac{1}{X^2} - \frac1X + \frac{1}{X-1} .
\]

$\dfrac{X+1}{(X-1)^3}$: $Y = X - 1$ के साथ अंश $Y + 2$ है:
\[
\frac{Y + 2}{Y^3} = \frac{1}{Y^2} + \frac{2}{Y^3}
= \frac{1}{(X-1)^2} + \frac{2}{(X-1)^3} .
\]
\end{solution}

\begin{solution}{exo:b1:fractions:4}
\omterm{def:b1:fractions:field}{ध्रुव} इकाई के चौथे मूल $1, \iu, -1, -\iu$ हैं, और सभी सरल हैं। $B' = 4X^3$
के साथ \omterm{def:b1:fractions:field}{सरल-ध्रुव} सूत्र: $a$ पर गुणांक $\frac{1}{4a^3} = \frac{a}{4a^4} =
\frac a4$ है ($a^4 = 1$ का प्रयोग करते हुए)। अतः $\C$ पर:
\[
\frac{1}{X^4 - 1}
= \frac{1/4}{X - 1} - \frac{1/4}{X + 1}
+ \frac{\iu/4}{X - \iu} - \frac{\iu/4}{X + \iu} .
\]
संयुग्मी जोड़े को समूहबद्ध करने पर (उभयनिष्ठ हर $X^2 + 1$):
$\frac{\iu}{4}\bigl(\frac{1}{X-\iu} - \frac{1}{X+\iu}\bigr) =
\frac{\iu}{4}\cdot\frac{2\iu}{X^2+1} = \frac{-1/2}{X^2+1}$। $\R$ पर:
\[
\frac{1}{X^4 - 1}
= \frac{1/4}{X-1} - \frac{1/4}{X+1} - \frac{1/2}{X^2 + 1} .
\]
\emph{$X = 0$ पर जाँच:} $-1 = -\frac14 - \frac14 - \frac12$।
\end{solution}

\begin{solution}{exo:b1:fractions:5}
$\dfrac{X^2}{(X^2+1)^2} = \dfrac{(X^2 + 1) - 1}{(X^2+1)^2} =
\dfrac{1}{X^2+1} - \dfrac{1}{(X^2+1)^2}$.

अतः एक प्रतिअवकलज:
\[
\int \frac{x^2\,\dd x}{(x^2+1)^2}
= \arctan x - \frac12\Bigl(\arctan x + \frac{x}{x^2+1}\Bigr) + C
= \frac12\arctan x - \frac{x}{2(x^2+1)} + C .
\]
\end{solution}

\begin{solution}{exo:b1:fractions:6}
\omterm{def:b1:fractions:field}{ध्रुव} $0, -1, \dots, -n$ सरल हैं। $-k$ पर आवरण विधि:
\[
c_k = \frac{n!}{\prod_{j \neq k} (-k + j)}
= \frac{n!}{\bigl(\prod_{j=0}^{k-1}(j - k)\bigr)
\bigl(\prod_{j=k+1}^{n}(j-k)\bigr)}
= \frac{n!}{(-1)^k k!\,(n-k)!} = (-1)^k \binom nk .
\]
अतः
\[
\frac{n!}{X(X+1)\cdots(X+n)}
= \sum_{k=0}^{n} \frac{(-1)^k \binom nk}{X + k} .
\]
($n = 1$ के लिए संगति की जाँच: $\frac{1}{X(X+1)} = \frac1X -
\frac{1}{X+1}$।)
\end{solution}

\begin{solution}{exo:b1:fractions:7}
$P = X^n - 1$ के $n$ सरल मूल $\omega^k$ (\cref{thm:b1:complex:roots}) हैं,
अतः \cref{exo:b1:poly:11} की लघुगणकीय-अवकलज सर्वसमिका यह कहती है
\[
\sum_{k=0}^{n-1} \frac{1}{X - \omega^k}
= \frac{P'(X)}{P(X)} = \frac{n X^{n-1}}{X^n - 1} .
\]
$X = 2$ पर $n = 4$: दायाँ पक्ष $= \frac{4 \times 8}{15} = \frac{32}{15}$।
बायाँ पक्ष: $\frac{1}{2-1} + \frac{1}{2+1} + \frac{1}{2 - \iu} + \frac{1}{2
+ \iu} = 1 + \frac13 + \frac{4}{5} = \frac{15 + 5 + 12}{15} =
\frac{32}{15}$, जहाँ $\frac{1}{2-\iu} + \frac{1}{2+\iu} = \frac{4}{5}$
प्रयुक्त हुआ।
\end{solution}

\begin{solution}{exo:b1:fractions:8}
आवरण विधि: $\dfrac{1}{k(k+1)(k+2)} = \dfrac{1/2}{k} - \dfrac{1}{k+1} +
\dfrac{1/2}{k+2}$। इसे दूरबीनी अंतर के रूप में फिर से लिखिए:
\[
\frac{1}{k(k+1)(k+2)}
= \frac12\Bigl(\frac{1}{k(k+1)} - \frac{1}{(k+1)(k+2)}\Bigr),
\]
(जाँचने के लिए प्रसार कीजिए --- अथवा दोनों वियोजनों को घटाइए)। योग करने पर:
\[
S_n = \frac12\Bigl(\frac{1}{1 \times 2} - \frac{1}{(n+1)(n+2)}\Bigr)
= \frac14 - \frac{1}{2(n+1)(n+2)}
\xrightarrow[n \to \infty]{} \frac14 .
\]
\end{solution}

\begin{solution}{exo:b1:fractions:9}
$0 \leq m \leq n - 1$ के लिए भिन्न $\frac{X^m}{P}$ (घात $m - n \leq -1$,
\omterm{def:b1:fractions:field}{ध्रुव} सरल) का वियोजन कीजिए: \omterm{def:b1:fractions:field}{सरल-ध्रुव} सूत्र देता है
\[
\frac{X^m}{P} = \sum_{i=1}^{n} \frac{x_i^m / P'(x_i)}{X - x_i} .
\]
$X$ से गुणा कीजिए और $X \to +\infty$ लीजिए: बायाँ पक्ष $X^{m+1}/P$ की सीमा
की ओर जाता है, जो $m \leq n - 2$ होने पर $0$ है और $m = n - 1$ होने पर $1$
($P$ घात $n$ का \omterm{def:b1:poly:def}{इकाई-अग्र} है); दायाँ पक्ष $\sum_i \frac{x_i^m}{P'(x_i)}$ की
ओर जाता है। अतः
\[
\sum_{i=1}^{n} \frac{x_i^{m}}{P'(x_i)} =
\begin{cases}
0 & \text{के लिए } 0 \leq m \leq n-2,\\
1 & \text{के लिए } m = n - 1,
\end{cases}
\]
जिसमें दोनों घोषित सर्वसमिकाएँ आ जाती हैं ($m = 0$ के लिए $n \geq 2$ आवश्यक
है)। (\cref{thm:b1:poly:lagrange} के द्वारा व्याख्या: ये योग $X^m$ के
लाग्रांज अंतर्वेशकों के अग्र गुणांक हैं, और घात $\leq n-1$ के \omterm{def:b1:poly:def}{बहुपद} का $n$
बिंदुओं पर अंतर्वेशन उसे ज्यों-का-त्यों लौटा देता है।)
\end{solution}

\begin{solution}{exo:b1:fractions:10}
आकार $\frac aX + \frac b{X+1} + \frac c{(X+1)^2}$। $0$ पर आवरण विधि: $a =
1$; द्विक \omterm{def:b1:fractions:field}{ध्रुव} पर आवरण विधि: $c = \bigl[\frac1X \bigr]_{X=-1} = -1$; अनंत
पर $xF(x)$ की सीमा: $0 = a + b$, अतः $b = -1$:
\[
\frac{1}{X(X+1)^2} = \frac1X - \frac1{X+1} - \frac1{(X+1)^2} .
\]
$n = 1, \dots, N$ के लिए योग करने पर: पहली दो ईंटें दूरबीनी होकर $1 -
\frac1{N+1}$ देती हैं, और तीसरी $-\sum_{k=2}^{N+1} \frac1{k^2}$ का योगदान
करती है। $N \to \infty$ लेने पर और $\sum_{k\geq1} \frac1{k^2} =
\frac{\pi^2}6$ का प्रयोग करने पर:
\[
\sum_{n=1}^{\infty}\frac{1}{n(n+1)^2}
= 1 - \Bigl(\frac{\pi^2}6 - 1\Bigr) = 2 - \frac{\pi^2}6
\approx 0.355 .
\]
\end{solution}

\begin{solution}{exo:b1:fractions:11}
चर $Y$ में: $\frac1{(Y+1)(Y+4)} = \frac{1/3}{Y+1} - \frac{1/3}{Y+4}$ ($-1$
और $-4$ पर आवरण विधि), अतः
\[
\frac{1}{(X^2+1)(X^2+4)}
= \frac13\,\frac1{X^2+1} - \frac13\,\frac1{X^2+4} .
\]
इसी प्रकार $\frac{Y}{(Y+1)(Y+4)} = \frac{-1/3}{Y+1} + \frac{4/3}{Y+4}$, अतः
\[
\frac{X^2}{(X^2+1)(X^2+4)}
= -\frac13\,\frac1{X^2+1} + \frac43\,\frac1{X^2+4} .
\]
($X = 0$ पर जाँच: $0 = \frac13(-1 + 1)$।) ये पहले से ही वास्तविक वियोजन हैं:
अखंडनीय द्विघातों के ऊपर के अंश संयोग से अचर निकल आते हैं।
\end{solution}

\begin{solution}{exo:b1:fractions:12}
\begin{enumerate}
  \item ध्रुवों से बचने वाले प्रत्येक अंतराल पर $F'(x) = -\sum_i
  \frac{c_i}{(x - p_i)^2} < 0$: अर्थात् कठोरतः ह्रासमान। जब $x \to
  \pm\infty$, तब प्रत्येक ईंट $0$ की ओर जाती है: $F \to 0$ --- $+\infty$ पर
  ऊपर से (वहाँ सभी ईंटें धनात्मक हैं) और $-\infty$ पर नीचे से। जब $x \to
  p_i^+$, तब ईंट $\frac{c_i}{x - p_i}$ $+\infty$ तक फट पड़ती है और शेष
  परिबद्ध रहती हैं: $F \to +\infty$; इसी प्रकार जब $x \to p_i^-$, तब $F \to
  -\infty$।
  \item $\lambda > 0$ नियत कीजिए। $\intoo{-\infty}{p_1}$ पर: $F$ $0^-$ से
  $-\infty$ तक घटता है, अतः $F < 0 < \lambda$: कोई हल नहीं। प्रत्येक
  $\intoo{p_i}{p_{i+1}}$ ($1 \leq i \leq r-1$) पर: $F$ $+\infty$ से
  $-\infty$ तक घटता है, अतः मान $\lambda$ ठीक एक बार लेता है (मध्यवर्ती मान
  गुणधर्म और कठोर एकदिष्टता से)। $\intoo{p_r}{+\infty}$ पर: $F$ $+\infty$ से
  $0^+$ तक घटता है, और फिर ठीक एक हल। कुल: ठीक $r$ हल, जो ध्रुवों के बीच-बीच
  में आते हैं।
\end{enumerate}
\end{solution}

\begin{solution}{pb:b1:fractions:1}
\textbf{1.} आवरण विधि: $\frac1{X(X+1)} = \frac1X - \frac1{X+1}$। योग दूरबीनी
हो जाता है:
\[
\sum_{n=1}^{N}\Bigl(\frac1n - \frac1{n+1}\Bigr)
= 1 - \frac1{N+1} \longrightarrow 1 .
\]

\textbf{2.} $\frac1{X(X+2)} = \frac{1/2}X - \frac{1/2}{X+2}$। योग करने पर $n
= 1, 2$ के लिए पद $\frac1n$ बच रहते हैं और $n = N-1, N$ के लिए पद
$-\frac1{n+2}$:
\[
\sum_{n=1}^{N}\frac1{n(n+2)}
= \frac12\Bigl(1 + \frac12 - \frac1{N+1} - \frac1{N+2}\Bigr)
\longrightarrow \frac34 .
\]

\textbf{3.} यदि $F(X) = G(X) - G(X+1)$, तो $\sum_{n=1}^N F(n) =
\sum_{n=1}^N\bigl(G(n) - G(n+1)\bigr) = G(1) - G(N+1)$: सभी मध्यवर्ती मान
जोड़ों में कट जाते हैं। $F = \frac1{X(X+1)(X+2)}$ के लिए साक्षी $G(X) =
\frac1{2X(X+1)}$ है:
\[
G(X) - G(X+1)
= \frac{(X+2) - X}{2X(X+1)(X+2)} = F(X) ,
\]
अतः $\sum_{n=1}^N F(n) = \frac14 - \frac1{2(N+1)(N+2)} \to \frac14$, जो
\cref{exo:b1:fractions:8} का मान है।

\textbf{4.} दाएँ पक्ष को उभयनिष्ठ हर $X(X+1)\cdots(X+k)$ पर रखिए:
\[
\frac1k\cdot\frac{(X + k) - X}{X(X+1)\cdots(X+k)}
= \frac{1}{X(X+1)\cdots(X+k)} ,
\]
जो वही सर्वसमिका है। अतः $G_k(X) = \frac1{X(X+1)\cdots(X+k-1)}$ के साथ
$F_k(X) = \frac1k\bigl(G_k(X) - G_k(X+1)\bigr)$, और प्रश्न 3 की क्रियाविधि
देती है
\[
\sum_{n=1}^{N}\frac1{n(n+1)\cdots(n+k)}
= \frac1k\Bigl(\frac1{k!} - G_k(N+1)\Bigr)
\longrightarrow \frac1{k\cdot k!} ,
\]
क्योंकि $G_k(1) = \frac1{k!}$ और $G_k(N+1) \to 0$। $k = 2$ के लिए:
$\frac1{2\cdot2} = \frac14$, जो प्रश्न 3 से मेल खाता है।

\textbf{5.} \cref{exo:b1:fractions:6} $\frac{n!}{X(X+1)\cdots(X+n)} =
\sum_{j=0}^n \frac{(-1)^j\binom nj}{X + j}$ देता है। $X = 1$ पर मान लीजिए:
बायाँ पक्ष $\frac{n!}{(n+1)!} = \frac1{n+1}$ है और दायाँ पक्ष
$\sum_j(-1)^j\binom nj\frac1{1+j}$: पहली सर्वसमिका। $X = 2$ पर: बायाँ पक्ष
$\frac{n!}{2\cdot3\cdots(n+2)} = \frac{n!\cdot1}{(n+2)!} =
\frac1{(n+1)(n+2)}$ है, दायाँ $\sum_j(-1)^j\binom nj\frac1{2+j}$: दूसरी
सर्वसमिका।

\textbf{6.} \omterm{def:b1:fractions:field}{ध्रुव} $\omega^k$ सरल हैं, और \cref{met:b1:fractions:compute} का
\omterm{def:b1:fractions:field}{सरल-ध्रुव} सूत्र गुणांक देता है
\[
\frac{1}{\bigl(nX^{n-1}\bigr)_{X = \omega^k}}
= \frac{1}{n\,\omega^{k(n-1)}}
= \frac{\omega^k}{n\,\omega^{kn}} = \frac{\omega^k}n ,
\]
जहाँ $\omega^{kn} = 1$ प्रयुक्त हुआ। अतः $\frac1{X^n-1} = \frac1n\sum_k
\frac{\omega^k}{X - \omega^k}$।

\textbf{7.} $n = 2$ ($\omega = -1$) के लिए: $\frac12\bigl( \frac1{X-1} -
\frac1{X+1}\bigr) = \frac12\cdot\frac{2}{X^2-1} = \frac1{X^2-1}$: सही। $n
\geq 2$ के लिए गुणांकों का योग $\frac1n\sum_k\omega^k = 0$ है
(\cref{prop:b1:complex:sumroots})। ऐसा होना ही चाहिए: सरल ध्रुवों वाले किसी
भी वियोजन के लिए जब $x \to \infty$ तब $x\,F(x) \to \sum_k c_k$, जबकि यहाँ
$xF(x) = \frac{x}{x^n-1} \to 0$, क्योंकि $n \geq 2$।

\textbf{8.} $\omega^k$ पर $\frac{X^{n-1}}{X^n - 1}$ के लिए आवरण विधि:
$\frac{A(\omega^k)}{B'(\omega^k)} = \frac{\omega^{k(n-1)}}{n\omega^{k(n-1)}}
= \frac1n$, अतः $\frac{X^{n-1}}{X^n-1} = \frac1n\sum_k\frac1{X - \omega^k}$
--- अर्थात् फिर से \cref{exo:b1:fractions:7} की सर्वसमिका।

\textbf{9.} $c = \omega^k$, $\conj c = \omega^{n-k}$ और $c\conj c = 1$, $c +
\conj c = 2\cos\theta_k$ के साथ:
\[
\frac{c}{X - c} + \frac{\conj c}{X - \conj c}
= \frac{c(X - \conj c) + \conj c(X - c)}
{(X - c)(X - \conj c)}
= \frac{2\cos\theta_k\,X - 2}{X^2 - 2\cos\theta_k\,X + 1} .
\]
प्रश्न 6 में $k$ को $n - k$ के साथ समूहबद्ध करने पर: विषम $n$ के लिए,
\[
\frac1{X^n - 1} = \frac1n\Biggl(\frac1{X-1} +
\sum_{k=1}^{(n-1)/2}
\frac{2\cos\theta_k X - 2}{X^2 - 2\cos\theta_k X + 1}\Biggr) ;
\]
सम $n$ के लिए अतिरिक्त स्वयं-युग्मित \omterm{def:b1:fractions:field}{ध्रुव} $\omega^{n/2} = -1$ कोष्ठक के
भीतर $\frac{-1}{X+1}$ का योगदान करता है और युग्म-योग $\frac n2 - 1$ तक चलता
है।

\textbf{10.} $n = 4$: $\theta_1 = \frac\pi2$, $\cos\theta_1 = 0$, अतः
युग्म-पद $\frac{-2}{X^2+1}$ है और
\[
\frac1{X^4-1} = \frac14\Bigl(\frac1{X-1} - \frac1{X+1}
- \frac2{X^2+1}\Bigr) ,
\]
जो \cref{exo:b1:fractions:4} का वियोजन है।

\textbf{11.} $P = \prod_{k=1}^{n-1}(X - \omega^k)$ ($X^n - 1$ को $X - 1$ से
भाग दीजिए), अतः \cref{exo:b1:poly:11} से $\frac{P'}P =
\sum_{k=1}^{n-1}\frac1{X-\omega^k}$। $X = 1$ पर: $P(1) = n$ और $P'(1) =
\sum_{j=1}^{n-1}j = \frac{n(n-1)}2$, जिससे
\[
\sum_{k=1}^{n-1}\frac1{1 - \omega^k}
= \frac{P'(1)}{P(1)} = \frac{n-1}2 .
\]

\textbf{12.} अर्ध-कोण: $1 - \eu^{\iu\theta} =
-2\iu\sin\frac\theta2\,\eu^{\iu\theta/2}$, अतः
\[
\frac1{1 - \eu^{\iu\theta}}
= \frac{\eu^{-\iu\theta/2}}{-2\iu\sin\frac\theta2}
= \frac{\iu\cos\frac\theta2 + \sin\frac\theta2}
{2\sin\frac\theta2}
= \frac12 + \frac\iu2\,
\frac{\cos\frac\theta2}{\sin\frac\theta2} .
\]
$\theta = \theta_k = \frac{2k\pi}n$ के साथ: $\frac1{1-\omega^k} = \frac12 +
\frac\iu2\cot\frac{k\pi}n$। $k = 1, \dots, n-1$ पर योग करके और प्रश्न 11 के
वास्तविक मान $\frac{n-1}2$ से तुलना करने पर: वास्तविक भाग ही सब कुछ बता देते
हैं, अतः $\sum_k\cot\frac{k\pi}n = 0$ --- जैसा सममिति $\cot\frac{(n-k) \pi}n
= -\cot\frac{k\pi}n$ भी दिखाती है।

\textbf{13.} $\frac{P'}P = \sum_k\frac1{X - \omega^k}$ का अवकलन करने पर:
\[
\frac{P''}P - \Bigl(\frac{P'}P\Bigr)^2
= -\sum_{k=1}^{n-1}\frac1{(X - \omega^k)^2} .
\]
$X = 1$ पर: $P''(1) = \sum_{j=2}^{n-1}j(j-1) = 2\binom n3 =
\frac{n(n-1)(n-2)}3$ (हॉकी-स्टिक सर्वसमिका, अथवा आगमन), अतः
\[
\sum_{k=1}^{n-1}\frac1{(1 - \omega^k)^2}
= \Bigl(\frac{n-1}2\Bigr)^2 - \frac{(n-1)(n-2)}3
= \frac{(n-1)\bigl(3(n-1) - 4(n-2)\bigr)}{12}
= \frac{(n-1)(5-n)}{12} .
\]

\textbf{14.} $c_k = \cot\frac{k\pi}n$ के साथ प्रश्न 12 के सूत्र का वर्ग करने
पर:
\[
\frac1{(1-\omega^k)^2}
= \Bigl(\frac12 + \frac\iu2 c_k\Bigr)^2
= \frac14 - \frac{c_k^2}4 + \frac\iu2\,c_k .
\]
योग करके और $\sum c_k = 0$ (प्रश्न 12) तथा प्रश्न 13 का प्रयोग करके:
$\frac{n-1}4 - \frac14\sum_k c_k^2 = \frac{(n-1)(5-n)}{12}$, अतः
\[
\sum_{k=1}^{n-1}\cot^2\frac{k\pi}n
= (n-1) - \frac{(n-1)(5-n)}3 = \frac{(n-1)(n-2)}3 .
\]
फिर $\frac1{\sin^2t} = 1 + \cot^2t$ $\sum_k\frac1{\sin^2\frac{k\pi}n} =
(n-1) + \frac{(n-1)(n-2)}3 = \frac{n^2-1}3$ देता है।

\textbf{15.} $n = 3$: $\cot^2\frac\pi3 + \cot^2\frac{2\pi}3 = \frac13 +
\frac13 = \frac23 = \frac{2\cdot1}3$; और $\frac1{\sin^2}$ का योग $\frac43 +
\frac43 = \frac83 = \frac{9-1}3$ है। $n = 4$: $1 + 0 + 1 = 2 =
\frac{3\cdot2}3$; और $2 + 1 + 2 = 5 = \frac{16-1}3$। दोनों सूत्र जाँच पर खरे
उतरते हैं।

\textbf{16.} $t \in \intoo0{\frac\pi2}$ के लिए चिरपरिचित तुलना $\sin t < t <
\tan t$ (क्षेत्रफल या उत्तलता का तर्क, जो उच्चतर माध्यमिक से परिचित है)
व्युत्क्रम लेने पर $\cot t < \frac1t < \frac1{\sin t}$ देती है, और वहाँ
तीनों धनात्मक हैं; वर्ग करने से क्रम बना रहता है। $t = \frac{k\pi}n$, $1
\leq k \leq m$, $n = 2m+1$ के साथ (अतः $t < \frac\pi2$):
\[
\cot^2\frac{k\pi}n < \frac{n^2}{k^2\pi^2} <
\frac1{\sin^2\frac{k\pi}n} .
\]

\textbf{17.} सममितियों $\cot^2\frac{(n-k)\pi}n = \cot^2\frac{k\pi}n$ और
$\sin^2$ के लिए भी वैसी ही सममितियों से प्रश्न 14 के योग आधे हो जाते हैं:
$\sum_{k=1}^{m}\cot^2\frac{k\pi}n = \frac{(n-1)(n-2)}6 = \frac{m(2m-1)}3$ और
$\sum_{k=1}^m\frac1{\sin^2\frac{k\pi}n} = \frac{n^2-1}6 = \frac{2m(m+1)}3$।
$k = 1, \dots, m$ पर प्रश्न 16 का योग करके और $\frac{\pi^2}{n^2}$ से गुणा
करके:
\[
\frac{\pi^2}{(2m+1)^2}\cdot\frac{m(2m-1)}3
\;<\; \sum_{k=1}^{m}\frac1{k^2} \;<\;
\frac{\pi^2}{(2m+1)^2}\cdot\frac{2m(m+1)}3 .
\]
जब $m \to \infty$, तब दोनों परिबंध $\frac{\pi^2}6$ की ओर जाते हैं (अनुपात
$\frac{m(2m-1)}{(2m+1)^2}$ और $\frac{2m(m+1)}{(2m+1)^2}$ दोनों $\frac12$ की
ओर जाते हैं), अतः दबाव से वर्धमान आंशिक योग अभिसरित होते हैं और
\[
\sum_{k=1}^{\infty}\frac1{k^2} = \frac{\pi^2}6 .
\]

\textbf{18.} $\frac1{X^2-1} = \frac12\bigl(\frac1{X-1} - \frac1{X+1}\bigr)$
का वर्ग कीजिए:
\[
\frac1{(X^2-1)^2} = \frac14\Bigl(\frac1{(X-1)^2} +
\frac1{(X+1)^2}\Bigr) - \frac12\cdot\frac1{(X-1)(X+1)} ,
\]
और आड़े पद $\frac1{(X-1)(X+1)} = \frac12\bigl(\frac1{X-1} -
\frac1{X+1}\bigr)$ का फिर से वियोजन करके बताया गया रूप पाइए। $X = 0$ पर:
बायाँ पक्ष $1$; दायाँ पक्ष $\frac14(1 + 1) - \frac14(-1 - 1) = \frac12 +
\frac12 = 1$।

\textbf{19.} $n \geq 2$ पर प्रश्न 18 का योग कीजिए। $S =
\sum_{k\geq1}\frac1{k^2} = \frac{\pi^2}6$ के साथ: $\sum_{n\geq2}
\frac1{(n-1)^2} = S$; $\sum_{n\geq2}\frac1{(n+1)^2} = S - 1 - \frac14$; और
दो-अंतराल वाली दूरबीन $\sum_{n\geq2}\bigl( \frac1{n-1} - \frac1{n+1}\bigr) =
1 + \frac12 = \frac32$। अतः
\[
\sum_{n=2}^{\infty}\frac1{(n^2-1)^2}
= \frac14\Bigl(2S - \frac54\Bigr) - \frac14\cdot\frac32
= \frac S2 - \frac{11}{16}
= \frac{\pi^2}{12} - \frac{11}{16} \approx 0.135 .
\]
संख्यात्मक रूप से: $\frac19 + \frac1{64} + \frac1{225} + \frac1{576} +
\frac1{1225} + \dots \approx 0.1111 + 0.0156 + 0.0044 + 0.0017 + 0.0008 +
\dots \approx 0.135$: संगत।

\textbf{20.} प्रश्न 8 की सर्वसमिका $\sum_k\frac1{X-\omega^k} =
\frac{nX^{n-1}}{X^n-1}$ का अवकलन कीजिए:
\[
\sum_{k=0}^{n-1}\frac1{(X - \omega^k)^2}
= \frac{n^2X^{2n-2} - n(n-1)X^{n-2}(X^n - 1)}{(X^n - 1)^2} .
\]
$n = 2$, $X = 2$ पर जाँच: दायाँ पक्ष $\frac{4\cdot4 - 2\cdot1\cdot3}{9} =
\frac{10}9$; बायाँ पक्ष $\frac1{(2-1)^2} + \frac1{(2+1)^2} = 1 + \frac19 =
\frac{10}9$।

\textbf{21.} $\frac1{\binom nk} = \frac{k!\,(n-k)!}{n!} =
\frac{k!}{(n-k+1)(n-k+2)\cdots n}$, जिसके हर में $k$ क्रमागत पूर्णांकों का
गुणनफल है। $j = n - k + 1$ प्रतिस्थापित करने पर (जिससे $n$ के $k$ से चलने पर
$j$ $\N^*$ पर चलता है):
\[
\sum_{n=k}^{\infty}\frac1{\binom nk}
= k!\sum_{j=1}^{\infty}\frac1{j(j+1)\cdots(j+k-1)}
= \frac{k!}{(k-1)\,(k-1)!} = \frac{k}{k-1} ,
\]
यह $k$ के स्थान पर $k - 1$ के साथ प्रश्न 4 लगाने से मिलता है (वैध, क्योंकि
$k - 1 \geq 1$)।

\textbf{22.} $k = 3$ के लिए: $\frac1{\binom n3} = \frac6{(n-2)(n-1)n}$, अतः
प्रश्न 3 के आंशिक योग (विस्थापित करके) से,
\[
\sum_{n=3}^{N}\frac1{\binom n3}
= 6\sum_{j=1}^{N-2}\frac1{j(j+1)(j+2)}
= 6\Bigl(\frac14 - \frac1{2(N-1)N}\Bigr)
= \frac32 - \frac3{(N-1)N}
\longrightarrow \frac32 ,
\]
अर्थात् प्रश्न 21 का मान $\frac k{k-1} = \frac32$।

\textbf{23.} $n = 6$: जोड़े $k = 1$ ($\theta_1 = \frac\pi3$, $2\cos\theta_1
= 1$) और $k = 2$ ($\theta_2 = \frac{2\pi}3$, $2\cos\theta_2 = -1$), साथ ही
वास्तविक \omterm{def:b1:fractions:field}{ध्रुव} $\pm1$:
\[
\frac1{X^6-1} = \frac16\Bigl(\frac1{X-1} - \frac1{X+1}
+ \frac{X - 2}{X^2 - X + 1}
+ \frac{-X - 2}{X^2 + X + 1}\Bigr) .
\]
$X = 0$ पर: बायाँ पक्ष $-1$; दायाँ पक्ष $\frac16(-1 - 1 - 2 - 2) = -1$: सही।

\textbf{24.} (क) अद्वितीयता गुणांकों की हर पहचान को वैध बनाती है --- आवरण
विधि, प्रश्न 9 का संयुग्मी-युग्म समूहन, और अवकलन की युक्तियाँ (प्रश्न 13,
20), सब उसी पर टिके हैं। (ख) इकाई के मूलों ने $X^n - 1$ के \omterm{def:b1:fractions:field}{ध्रुव}, उनकी
सममितियाँ ($k \leftrightarrow n-k$) और प्रश्न 12 का अर्ध-कोण बीजगणित
(\cref{met:b1:complex:trig}) दिया। (ग) लघुगणकीय अवकलज $\frac{P'}P =
\sum\frac{m_i}{X-a_i}$ ने $P = 1 + X + \dots + X^{n-1}$ के मूलों की सूचना को
प्रश्न 11 और 13 के संख्यात्मक योगों में बदल दिया --- अर्थात् भाग II और III
के बीच की कब्ज़ा।

\textbf{25.} वियोजन विशुद्ध रूप से बीजीय सर्वसमिका है, जो चर के हर मान के
लिए एक साथ सत्य है; यही उसे इंजन बनाता है। पूर्णांक प्रतिस्थापित करके योग
करने पर वह दूरबीन बन गया (भाग I); \omterm{def:b1:complex:unity}{इकाई के मूल} प्रतिस्थापित करके और संयुग्मी
समूहबद्ध करके वह त्रिकोणमितीय सर्वसमिकाएँ बन गया (भाग II और III); और
विश्लेषण केवल सबसे अंतिम पद पर आया --- दो संवृत रूपों के बीच का दबाव ---
जिससे $\frac{\pi^2}6$ मिला, ऐसा \omterm{def:b1:logic:statement}{कथन} जिस तक कोई परिमित प्रतिस्थापन नहीं पहुँच
सकता था। यह श्रम-विभाजन (बीजगणित यथार्थ परिमित सर्वसमिकाएँ देता है, विश्लेषण
सीमा तक ले जाता है) \cref{ch:b1:series} का आदर्श है, जहाँ दूरबीनी और तुलना
व्यवस्थित हो जाती हैं, और \cref{ch:b1:integration} का भी, जहाँ हर ईंट अपना
प्रतिअवकलज पा लेती है और वही वियोजन योगों के बदले समाकल संगणित करते हैं।
\end{solution}
